%% kapa: exam=1 round=30 year=2019 period=2 form=A subject=accounting subject_ko=회계학 seq=34 number=74
%% source: Q-Net 감정평가사 기출문제 (cst00309) artlSeq=5205205 file=extracted/1008/5205205/2019년 제30회 감정평가사 1차시험 문제지/2019년 제30회 감정평가사 1차시험 문제지/2019년도 제30회 감정평가사 1차 2교시 A형.hwp
%% conversion: kordoc 4.24.51 -> tools/gen1.py (auto); review=proofread-sonnet; flags=none
\begin{question}{74}
(주)감평은 활동기준원가계산방법에 의하여 제품의 원가를 계산하고 있다. 다음은 (주)감평의 연간 활동제조간접원가 예산자료와 작업 \#203의 원가동인에 관한 자료이다.
\begin{framedtext}
○ 연간 활동제조간접원가 예산자료 \par
{\small\centering
\begin{tabular}{|c|c|c|c|}
\hline
활동 & 활동별 제조간접원가 & 원가동인 & 원가동인수량 \\
\hline
생산준비 & ￦200,000 & 생산준비시간 & 1,250시간 \\
\hline
재료처리 & ￦300,000 & 재료처리횟수 & 1,000회 \\
\hline
기계작업 & ￦500,000 & 기계작업시간 & 50,000시간 \\
\hline
품질관리 & ￦400,000 & 품질관리횟수 & 10,000회 \\
\hline
\end{tabular}\par}
\vspace{0.5em}
○ 작업 \#203의 원가동인 자료 \par
{\small\centering
\begin{tabular}{|c|c|c|c|c|}
\hline
작업 & 생산준비시간 & 재료처리횟수 & 기계작업시간 & 품질관리횟수 \\
\hline
\#203 & 60시간 & 50회 & 4,500시간 & 500회 \\
\hline
\end{tabular}\par}
\end{framedtext}
작업 \#203의 제조원가가 ￦300,000이라면, 작업 \#203의 기본(기초)원가는? \par
\begin{choices}
  \item ￦210,400
  \item ￦220,000
  \item ￦225,400
  \item ￦230,400
  \item ￦255,400
\end{choices}
\end{question}
\answer{1}
